\documentclass[../main.tex]{subfiles}
\begin{document}

\FloatBarrier
\section{Problem Statement}
\label{sec:problem}

Online and hybrid education is now a permanent part of teaching, and lectures, tutorials, and exams are given through video calls in which the webcam is the main link between student and teacher. The webcam stream shows how well a student follows the material, but a teacher facing dozens of simultaneous video tiles cannot notice when one learner loses focus. Automatic engagement recognition addresses this problem: it classifies a short webcam clip of one student into one of four ordinal levels, from E0 (Highly\nobreakdash-Disengage) to E3 (Highly\nobreakdash-Engage), so a teacher or platform can adapt the pace of a lesson.

This thesis adopts the task form used by prior work and by the datasets it builds on. The input is a short front-webcam clip of one student, about ten seconds long. The output is one of four ordinal levels, E0 (Highly\nobreakdash-Disengage), E1 (Disengage), E2 (Engage), and E3 (Highly\nobreakdash-Engage), a scale from full withdrawal to active attention.

Although it resembles ordinary video classification, engagement recognition is harder for three reasons that shape every design choice in this thesis.

\begin{enumerate}
  \item \textbf{Ordinal labels.}
  The four classes lie on a line, \mbox{E0 $<$ E1 $<$ E2 $<$ E3}, so predicting E3 for an E0 student is a serious error while predicting E2 for an E3 student is almost right. Accuracy ignores this order and treats both mistakes as equally wrong, so the size of an error must be weighted by the distance between the true and predicted levels.

  \item \textbf{Severe imbalance.}
  Real classes are mostly attentive, so the datasets are dominated by the upper levels: about 69\% of CMOSE clips are E2 and fewer than 3\% are E0 (Section~\ref{sec:data}, Table~\ref{tab:T3_1_dataset_stats}). A model can therefore reach high accuracy by always predicting the majority class and never flagging the rare disengaged cases a monitoring system must catch.

  \item \textbf{Unfixed camera angle.}
  Engagement is read from facial behaviour, but the camera geometry differs from learner to learner: a laptop camera looks up from the desk, a monitor camera sits near eye level, a phone is propped to the side. The OpenFace gaze, head-pose, and landmark features are measured in image geometry, so a change of angle moves the features even when the underlying engagement is identical. Each dataset covers its own narrow range of setups, so a deployed model meets unfamiliar viewpoints immediately, yet robustness to this shift is rarely measured.
\end{enumerate}

These three properties follow from how an engagement signal is used, and together they define what a useful model must do. The unit of decision is rarely a single clip but a summary, such as an engagement curve for one student over a session or a class-level spread the teacher watches live. The value of that signal lies in the rare cases: confirming that an attentive class is attentive tells the teacher little, while flagging the few moments a student disengages is the goal, and this is the case that imbalance hides and accuracy does not reward. A deployment also never sees the data the model was tuned on, since cameras, lighting, participants, and rubric all change. The main concern of this thesis is therefore not overall accuracy but ordinal agreement on imbalanced data, and how well it holds when the data shifts, judged on a self-collected, never-trained-on test set under a cross-dataset protocol.

\FloatBarrier
\section{Background and Problems of Research}
\label{sec:background}

Vision-based engagement recognition has taken two routes. The first encodes each frame with a facial-behaviour toolkit, most often OpenFace~\cite{openface2}, which estimates gaze direction, head pose, facial landmarks, and action-unit intensities, and feeds the feature sequence to a classical classifier or a recurrent network such as an OpenFace$+$BiLSTM pipeline~\cite{openface_bilstm} or classical models on summary features~\cite{classical_ml_engagement}. The second learns directly from pixels with deep video models, for instance an EfficientNetV2 backbone followed by an LSTM~\cite{efficientnet_lstm_engagement}, or borrows context-modelling ideas from affect recognition~\cite{context_aware_emotion_3d}. This thesis builds on two in-the-wild datasets: CMOSE~\cite{cmose}, a recent, large, expert-labelled dataset that ships ready OpenFace and I3D features, and DAiSEE~\cite{daisee}, a widely used public dataset that gives only raw video, from which the same features are extracted. A systematic review~\cite{review_engagement_cv} confirms that class imbalance and subjective labels are common, unsolved problems across the field.

Two limits of this work motivate the thesis. The first is that the mixed facial-behaviour features are passed through a single encoder. A single OpenFace frame in CMOSE is a 709-dimensional vector that bundles signals of different physical kinds: a few gaze values, hundreds of eye- and face-landmark coordinates, head-pose parameters, and action-unit strengths. These parts move at different speeds, since the eyes flick fast and noisily, head turns are slow and smooth, and action units fire in short bursts. Joining all 709 numbers and pushing them through one sequence encoder forces a single temporal model onto signals that do not share a timescale, and hides the role of each behavioural channel.

The second limit is that studies rely on accuracy as the primary metric and lack a reliable evaluation of generalisation. Almost every study reports accuracy in-domain, training and testing on the same dataset, yet the only deployment that matters is on data recorded with different cameras and conditions. Because engagement labels are subjective and the class balance differs sharply between datasets, models can be expected to transfer poorly, but without measuring transfer this stays an assumption. On imbalanced data a model that predicts the majority class can reach a decent accuracy while learning nothing that transfers, so the metric most papers report hides the failure it should reveal.

These two limits define the gap. The mixed-feature problem calls for a representation that keeps the structure of the facial features instead of flattening it, and the generalisation problem calls for a protocol that measures transfer directly with metrics that majority prediction cannot fool.

\FloatBarrier
\section{Research Objectives and Conceptual Framework}
\label{sec:objectives}

The goal of this thesis is to improve, and describe accurately, how engagement is recognised on a four-level ordinal scale from facial-behaviour and motion features. This thesis pursues it through one architecture idea and one set of evaluation rules, stated as four research questions.

\begin{enumerate}
  \item \textbf{RQ1: Decomposition.}
  Does splitting the OpenFace features into groups, each with its own temporal encoder, improve engagement classification over single-encoder baselines that encode all 709 features together?
  \item \textbf{RQ2: Encoder assignment.}
  If splitting helps, which group$\rightarrow$encoder choices matter? Does any behavioural channel prefer one encoder family (convolutional, recurrent, or attention-based)?
  \item \textbf{RQ3: Motion fusion.}
  Does adding a global I3D motion stream help, and does any gain hold across many configurations rather than one favourable one?
  \item \textbf{RQ4: Generalisation and losses.}
  How well do these models generalise across datasets, and how does the loss trade ordinal agreement against minority-class recall and balanced ordinal error?
\end{enumerate}

The idea behind RQ1--RQ3 is to treat a clip as several streams in time rather than one sequence, realised by the proposed hybrid model of Chapter~3 (Figure~\ref{fig:hybrid}). The OpenFace features are split into five feature groups (gaze, eye landmarks, face landmarks, head pose, and action units), and each group is encoded on its own into a fixed-length embedding by an encoder suited to its timing. A global I3D clip embedding can be added as a sixth stream. The embeddings are joined and classified, and every stream also gets its own small head so each channel is useful on its own. This design is the proposed hybrid model, the main object tested against the single-encoder baselines. The rule behind RQ4 is to select and rank every model on quadratic-weighted kappa (QWK), an ordinal-agreement metric, reporting balanced accuracy and balanced ordinal error (macro-MAE) only as secondary checks, and to test every model not only in-domain but on three datasets, one collected for this thesis.

\FloatBarrier
\section{Contributions}
\label{sec:contributions}

This thesis makes three contributions.

\begin{enumerate}
  \item \textbf{Architecture.}
  This thesis proposes a hybrid model that splits the 709 OpenFace features into five groups, encodes each group with its own TCN, Transformer, or LSTM, optionally adds a global I3D motion stream, and trains every stream with an auxiliary head (Section~\ref{sec:hybrid}). Trying all 243 encoder assignments, with the I3D stream enabled and disabled, shows the I3D-stream-enabled configurations score higher than the strongest baseline (median QWK 0.553 against 0.537, best 0.605), that the gain belongs to the family rather than one configuration, and that head pose is the only group with a clear encoder preference, for a TCN, which holds under shift (Section~\ref{sec:ablation}).

  \item \textbf{Dataset.}
  A set of 366 webcam clips was collected, cut to a fixed length, filtered for tracking quality, labelled by hand, and processed with the same OpenFace/I3D pipeline as the public datasets. This test-only set measures performance on real students recorded on their own webcams, where camera angles and conditions differ from the public datasets (Section~\ref{sec:data}).

  \item \textbf{Generalisation.}
  This thesis runs a 3$\times$3 cross-dataset study (training on CMOSE, DAiSEE, or their combination; testing on each public test set plus the private set) showing that engagement models do not transfer to a new dataset, that accuracy hides this collapse while QWK exposes it, and that combined training is the simplest effective fix. On the private set the architecture and combined training reinforce each other: the proposed hybrid model trained on both datasets is the best model (QWK 0.379 against 0.285 for the best baseline, a wider gap than in-domain), which shows the value of the three contributions together (Sections~\ref{sec:crossdomain} and~\ref{sec:private}).
\end{enumerate}

A reproducible pipeline turns about 1{,}500 trained runs into the figures and tables of this thesis, from feature extraction through training, cross-dataset evaluation, and figure and table generation.

\FloatBarrier
\section{Organization of Thesis}
\label{sec:organization}

The remainder of this thesis is organised as follows.

Chapter~\ref{chap:litreview} reviews vision-based engagement recognition and the background the thesis relies on. It sets the scope, compares the feature-based and end-to-end approaches with their limits, and then presents the background knowledge later chapters assume: the OpenFace and I3D feature extractors, the three temporal architecture families (TCN, LSTM, Transformer), the loss functions for imbalanced ordinal labels, the ordinal metrics, and the statistical measures used to analyse results.

Chapter~\ref{chap:method} presents the methodology. It gives an overview of the pipeline with diagrams of the baseline and the proposed model, describes the three training datasets and their imbalance, the private test set, and the leak-free preprocessing, defines the single-encoder baselines, and builds the feature-group decomposition and the proposed hybrid model on top of them.

Chapter~\ref{chap:results} reports the experimental results, centred on the proposed hybrid model. It sets out the loss forms, the training regime, and the 3$\times$3 cross-dataset evaluation matrix; confirms that the evaluation protocol holds across the model's 486 configurations; answers the design questions (which encoder suits which group, and whether the I3D stream helps); and takes the model out of domain, where it outperforms the single-encoder baseline in every cell of the matrix and, with combined training, reaches the best result on the private set (RQ4).

Chapter~\ref{chap:conclusions} draws the findings together, restates the contributions in the light of the evidence, discusses the limits, and proposes future work: targeting the rarest disengaged class, explicit domain adaptation, and learned group encoders.

Appendix~\ref{chap:baseline} presents the supporting single-encoder baseline study, protocol-first: it establishes the metrics (QWK, macro-accuracy, macro-MAE), the development dataset (CMOSE), and the development loss (cross-entropy), reads the in-domain leaderboard, studies prediction agreement to show the OpenFace and I3D families make different errors, and takes the baselines out of domain in the 3$\times$3 matrix, setting the reference level the proposed model is measured against.

\end{document}
